\section{Background}
\label{sec:background}

\subsection{UAV Communication Networks}

A Flying Ad-hoc Network (FANET) is a specific class of Mobile Ad-hoc Network (MANET) composed of small, autonomous UAVs that communicate with each other and with ground control stations in a decentralized manner. Unlike terrestrial MANETs, FANETs are characterized by high node mobility, frequent topology changes, and three-dimensional movement, which introduce significant challenges for routing, quality of service (QoS) provisioning, and network stability \cite{hafeez2023blockchain}. The communication links in FANETs are often intermittent and subject to interference, necessitating robust and adaptive protocols that can maintain connectivity in dynamic and often harsh environments. Key research areas in UAV communications include the development of efficient routing protocols, medium access control (MAC) strategies, and security mechanisms tailored to the unique constraints of aerial networks \cite{luo2023escm}.

\subsection{Blockchain Technology}

Blockchain is a distributed, immutable ledger technology that first gained prominence as the foundation of cryptocurrencies like Bitcoin. At its core, a blockchain is a chain of blocks, where each block contains a batch of transactions, a timestamp, and a cryptographic hash of the previous block, creating a tamper-resistant and chronologically ordered record of events \cite{jagatheesaperumal2023blockchain}. The ledger is maintained across a peer-to-peer network, and new blocks are added through a consensus mechanism, a protocol that allows distributed nodes to agree on the state of the ledger without a central coordinator. Common consensus algorithms include Proof-of-Work (PoW), Proof-of-Stake (PoS), and Practical Byzantine Fault Tolerance (PBFT). The key attributes of blockchain—decentralization, transparency, immutability, and security—make it a compelling technology for applications beyond finance, including supply chain management, healthcare, and, increasingly, secure communication systems \cite{hossain2024blockchain}.
